\section{Length-sensitive transport of actual clearance sources}
\label{sec:lc-length}

We keep the physical model, radius interval, original section and unit-speed
clock fixed.  In this section a marked closest point is a coordinate on a
free line, never a reflection.  The purpose is to retain the length of a
small source in a complete branch-supremum estimate.  An unscaled source
mass or a bound at one common endpoint momentum would not supply that
estimate.  The full-word pressure theorem of Section~\ref{sec:bp-pressure}
is not used below.

Use the enlarged physical curve class of \eqref{eq:mc-class}, with its own
curvature allowance $\widehat B_{\rm cv}$.  Set
\[
 t_*=3/50,\quad k_-=100/47,\quad k_+=20/9,\quad U=170/9,
 \quad \sigma=47/59,\quad \Gamma=\sqrt{1+k_+^2}.
\]
Choose its own homogeneous cutoff $\widehat k_0$ so that
\begin{equation}\label{eq:lc-theta}
 \theta_{\rm r}:=\sigma+
       \frac{J_f}{k_-t_*(\widehat k_0-1)}
                 <\frac{1+\sigma}{2}<1.
\end{equation}
The one-step and distortion proofs in
Theorems~\ref{thm:uc-expansion} and \ref{thm:uc-distortion} apply with
this fixed allowance and cutoff.  Take a common short length $\ell>0$
small enough for those proofs, the first-hit geometry, and the
relative-angle condition below.  The choices are independent of radius,
collision count, mark level, original exact labels and outer parameter $B$.
Let $D$ be the resulting Euclidean distortion constant and put
\begin{equation}\label{eq:lc-renewal-constants}
 D_\ell=e^{D\ell^{1/3}},\qquad
 C_\ell=\frac{3\Gamma D_\ell}{\ell(1-\theta_{\rm r})},\qquad
 G_\ell=\Gamma+C_\ell\ell.
\end{equation}

\subsection{The complete renewal, including the last generation}

At each step take every maximal regular physical homogeneous image.
Only images longer than $\ell$ are subdivided, into pieces of lengths
in $[\ell/2,\ell]$.  Images of length at most $\ell$ are not artificially
subdivided.  Write $\mathcal G_n(C)$ for the complete resulting partition
of the $n$th image of a qualified curve $C$, retaining each source label.
Call a piece long when its length is at least $\ell/3$.

\begin{theorem}[Length-sensitive physical renewal]\label{thm:lc-renewal}
For every qualified homogeneous curve $C$ of length at most $\ell$,
every $n\ge0$ and every original radius,
\begin{equation}\label{eq:lc-renewal}
 \sum_{V\in\mathcal G_n(C)}
        \sup_V J_V(T_R^{-n})
       \le \Gamma\theta_{\rm r}^{,n}+C_\ell|C|
       \le G_\ell.
\end{equation}
The sum contains all regular physical descendants; the omitted singular
source points have arclength zero.  The same statement holds after
physical reversal.  It is an estimate of actual inverse Jacobian
suprema, not an average or an independent-branch model.
\end{theorem}
\begin{proof}
First consider descendants having no long piece at the positive
generations $1,\ldots,n$.  They cannot pass through a length subdivision,
since every child of such a subdivision has length at least $\ell/2$.
At each generation their inverse adapted-Jacobian sum is therefore
bounded by the full maximal homogeneous sum in
\eqref{eq:uc-onestep}.  Products and nonnegative summation give
$\theta_{\rm r}^n$.  The endpoint equivalence of the adapted and
Euclidean curve norms costs $\Gamma$ once for the entire $n$-step
Jacobian, and not once per step.  Their total contribution is at most
$\Gamma\theta_{\rm r}^n$.  This also covers $n=0$.

Every other terminal piece has a unique last long ancestor $W$ at a
time $k\in\{1,\ldots,n\}$.  Denote its original preimage by $C_W$.
Its images are regular and homogeneous, and $\ell/3\le |W|\le\ell$.
The Euclidean distortion theorem yields
\[
 \sup_W J_W(T_R^{-k})
 \le D_\ell\frac{|C_W|}{|W|}
 \le \frac{3D_\ell}{\ell}|C_W|.
\]
The descendants after this last long ancestor have no later long
piece.  The preceding adapted-norm argument bounds the complete sum
of their inverse suffix Jacobians by
$\Gamma\theta_{\rm r}^{n-k}$.  For every fixed $k$, the genuine
preimage intervals $C_W$ of the long pieces are disjoint in $C$;
therefore $\sum_W|C_W|\le |C|$.  The contribution of all last-long
ancestors is consequently at most
\[
 \frac{3\Gamma D_\ell}{\ell}|C|
       \sum_{k=1}^{n}\theta_{\rm r}^{n-k}
 \le C_\ell|C|.
\]
The endpoint $k=n$ is included and has exponent zero.  No terminal
long piece is discarded.  All calculations can first be made for a
finite subfamily; nonnegative exhaustion gives the complete sum.
The physical one-step singularities and homogeneous boundaries have
countably many preimages on the qualified curves, with zero arclength
as in the native growth proof.  Reversal is the exact arclength
isometry conjugating $T_R$ to $T_R^{-1}$; it transports the original
observations and does not replace the section.  This proves
\eqref{eq:lc-renewal}.
\end{proof}

\subsection{A short interval cover of the true clearance event}

Write $N_*$ for the complete candidate-center bound of
\eqref{eq:ag-packing}, using the proved physical horizon, and put
\begin{equation}\label{eq:lc-cover-constants}
 \delta_*=3/200,\qquad a_{\rm c}=9/200,\qquad
 K_{\rm c}=2N_*,\qquad
 L_{\rm c}=\frac{2N_*\sqrt{1+U^2}}{k_-a_{\rm c}}.
\end{equation}
Let $E_{\delta,R}=\{N_{\delta,R}\ge1\}$ be the original event that
the actual outgoing flight has an unvisited clearance witness of
width $\delta$.  All physical visibility and interior-foot conditions
in that definition are retained.

\begin{lemma}[Actual-flight clearance cover]\label{lem:lc-cover}
For $0<\delta\le\delta_*$, on any qualified short graph $C$ the
set $C\cap E_{\delta,R}$ is covered, apart from endpoints, by at
most $K_{\rm c}$ disjoint source intervals whose combined arclength
is at most $L_{\rm c}\delta$.  The same bound holds on a short
horizontal momentum graph, using $0\le V\le U$ in the proof.
The cover consists of true outgoing states.  It can be larger than
the selected event, but every subsequent evolution always takes its
actual first hit.
\end{lemma}
\begin{proof}
Use boundary arclength $s$, outgoing angle $\varphi$, physical
velocity $v$ and $V=d\varphi/ds$.  For a candidate center $z$ put
\[
 t_z=(z-q)\cdot v,\qquad b_z=(z-q)\cdot v^\perp,
 \qquad \eta_z=\arg(z-q)-\arg v.
\]
Differentiating the actual circular geometry gives
\begin{equation}\label{eq:lc-impact-derivatives}
 b_z'=-\cos\varphi-t_z(\kappa+V),\qquad
 t_z'=-\sin\varphi+b_z(\kappa+V),\qquad \kappa=1/R.
\end{equation}
Since $|z-q|\ge1-R_+$, we have
$\eta_z'\le(1-R_+)^{-1}-k_-<0$ even when $V=0$.
The short length was chosen so that the variation of each $\eta_z$
is at most $\pi/2$.  Thus its positive-foot set is an interval.
On that interval $b_z$ is strictly decreasing.

The two bands $R<|b_z|<R+\delta$, with $t_z>0$, each give at most
one interval.  Throughout either band,
\[
 t_z^2=|z-q|^2-b_z^2
 \ge(1-R_+)^2-(R_++\delta_*)^2
 =\frac{1827}{40000}>a_{\rm c}^2.
\]
Equation~\eqref{eq:lc-impact-derivatives} therefore gives
$|b_z'|\ge k_-a_{\rm c}$.  The length in $s$ of each band is at
most $\delta/(k_-a_{\rm c})$.  Conversion to curve arclength costs
at most $\sqrt{1+U^2}$.  Sum over both signs and the at most $N_*$
possible centers, and merge overlapping intervals.  Merging decreases
both the number of components and their combined length.

Every genuine witness belongs to this cover.  We did not impose the
upper foot restriction $t_z<\tau$ when forming the cover.  A covered
state whose candidate foot is beyond an earlier collision is still
evolved to that earlier first collision; its inclusion only enlarges
a positive source estimate.  Endpoints of the weak bands do not
alter the length bound.  On a regular graph each such endpoint is
isolated because of the strict derivative just proved.  The same
reasoning, with $V=0$, proves the horizontal assertion.
\end{proof}

\subsection{The two physical source classes and their exact weights}

We use (i) the two marked angular circles of
Theorem~\ref{thm:ag-family}, with source coordinate $\xi=\vartheta$
and measure $\sum_{\varsigma=\pm1}d\vartheta$, and (ii) the horizontal
momentum fibre $F_a$, $a\in(-1,1)$, with $\xi=\alpha_0$ and measure
$d\alpha_0$.  Their respective masses are $4\pi$ and $2\pi$.
For (i), $u_0\in[0,\delta_*]$; at $u_0=0$ use the original strict
positive-side unvisited trace, with no reflection at the foot.
Both first-image families are now formed with the same enlarged
class, cutoff $\widehat k_0$ and length $\ell$ fixed above.

Here are fully specified uniform source budgets.  Set
\begin{align}
 M_*&=8N_*+1,&
 L_0&=\frac{\pi\sqrt{1+U^2}}{k_-},\notag\\
 J_0&=2+\left\lceil\frac{2\pi R_+}{\ell}\right\rceil,&
 J_{\rm len}&=2+\left\lceil
   \frac{2\max\{L_0,2\pi R_+\sqrt{1+U^2}\}}{\ell}\right\rceil,
                                                    \label{eq:lc-source-counts}\\
 A_{\rm ang}&=\frac{2M_*J_{\rm len}}{a_\circ}
                  \left(1+\frac2{\widehat k_0-1}\right),&
 A_{\rm hor}&=\frac{J_0J_fJ_{\rm len}}{t_*}
                  \left(1+\frac2{\widehat k_0-1}\right),\notag\\
 A_{\rm src}&=\max\{A_{\rm ang},A_{\rm hor}\},&
 E_{\rm lc}&=\frac1{k_-t_*}=\frac{47}{6},\notag\\
 H_{\rm lc}&=\max\left\{2M_*E_{\rm lc},J_0J_fE_{\rm lc},
                  \frac{A_{\rm src}G_\ell J_f}{k_-t_*}\right\}.
                                                    \label{eq:lc-source-budgets}
\end{align}
Shorten $\ell$ also below $\ell_{\rm pre}$ in \eqref{eq:mc-pre}.
The exact first-image formulas \eqref{eq:ag-density} and
\eqref{eq:mc-density} then show that the sum of density suprema in
either family is at most $A_{\rm src}$.  These bounds retain the
initial grazing factor $c$ before summing every homogeneous strip.
They are valid at $u_0=0$ by the direct trace construction of
Theorem~\ref{thm:nm-trace}.  The bounds are not conditional
probabilities on the source fibres.

\begin{theorem}[Localized half-source capacity]\label{thm:lc-local-half}
Start from either of the preceding physical sources and take its first
true collision $Y$.  Fix $h,n\ge0$ and select
$T_R^hY\in E_{\delta,R}$.  Propagate $n$ further true collisions
from that selected outgoing state, and then one last true collision.
Thus the half-count from the original source is exactly
\begin{equation}\label{eq:lc-half-index}
                         l=h+n+2.
\end{equation}
Let $A_l=|\partial_\xi p_l|$ be the actual final momentum derivative.
There is a countable disjoint family of regular source intervals,
with fixed physical half-words, covering the selected event up to
source-null sets, such that
\begin{equation}\label{eq:lc-local-half}
 \sum_I\sup_I A_l^{-1}
     \le K_{\rm loc}(\delta+\theta_{\rm r}^{,n}),\qquad
 K_{\rm loc}:=
 \frac{A_{\rm src}G_\ell J_f}{k_-t_*}
                  (\Gamma K_{\rm c}+C_\ell L_{\rm c}).
\end{equation}
Every individual interval has the bound
\begin{equation}\label{eq:lc-individual}
                \sup_I A_l^{-1}\le E_{\rm lc}\sigma^{l-1}.
\end{equation}
The complete unselected half-source, at every $l\ge1$, has a
corresponding interval partition with sum at most $H_{\rm lc}$ and
the same individual bound.  All assertions hold in the reversed
physical direction with the original observations transported.
\end{theorem}
\begin{proof}
Push the first-image family through $h$ true collisions.  By
Theorem~\ref{thm:lc-renewal}, the sum of the suprema of all exact
transported densities on the resulting short homogeneous curves
is at most $A_{\rm src}G_\ell$.  On each such curve use the disjoint
clearance cover of Lemma~\ref{lem:lc-cover}.  Apply the renewal
inequality separately to its at most $K_{\rm c}$ intervals through
$n$ further true collisions.  Their combined inverse-Jacobian sum
is at most
\[
                 \Gamma K_{\rm c}\theta_{\rm r}^n
                            +C_\ell L_{\rm c}\delta.
\]
This is where both component count and actual selected length are
paid.  A small mass is not substituted for an inverse-Jacobian
supremum.

On each resulting preterminal curve use the exact last-collision
cancellation \eqref{eq:mc-last-cancellation}.  The derivative of
final momentum per incoming arclength is at least $k_-t_*$, and
there are at most $J_f$ genuine next-hit pieces.  Do not make an
infinite terminal homogeneous subdivision.  Chain the exact source
density, inverse transport Jacobian and terminal derivative; then
sum over the last-hit cuts, selected covers, transported curves and
first-image family.  This gives the first inequality in
\eqref{eq:lc-local-half}.  All intervals are pullbacks through
injective physical branches, so disjoint covers remain disjoint
in the original source.  Intermediate physical and homogeneous
cuts have the native zero-arclength preimages; a bounded Borel
observation can be retained inside the resulting positive integrals
without differentiating it or counting its components.

For the individual inequality use the actual positive transverse
Jacobi product, not the averaged density.  On a marked angular
source its first incoming-side entries are $B=a^{\rm foot}\ge
 a_\circ>t_*$ and $D=1$.  On a horizontal momentum source they
are $B=R\cos\varphi_0+\tau\ge t_*$ and $D=1$; its exact first
endpoint derivative is $\cos\varphi_0+\tau/R+c_1$.
Each intervening kick and following physical flight replaces them by
\[
 B_{\rm new}=(1+\tau k)B+\tau D,\qquad
 D_{\rm new}=kB+D,\qquad k=2/(Rc)>0.
\]
Since $\tau\ge t_*$ and $c\le1$, the factor of $B$ is at least
$1+2t_*/R_+=59/47$.  The terminal endpoint is incoming: its
momentum derivative is $\kappa B+c_lD$, with no additional kick.
It is therefore at least $k_-t_*\sigma^{-(l-1)}$.
This proves \eqref{eq:lc-individual} pointwise and on every
subinterval, with no positive lower incidence cutoff.

For the complete unselected source with $l\ge2$, use the first
family, renewal through $l-2$ collisions and the same final
cancellation.  The sum is at most
$A_{\rm src}G_\ell J_f/(k_-t_*)$.  At $l=1$ the marked circles
have at most $2M_*$ first-hit/seam intervals; the horizontal fibre
has at most $J_0J_f$.  Apply the individual bound.  These are
exactly the three terms in $H_{\rm lc}$.  Direct construction at
the strict zero-clearance trace uses the same first-image density,
standoff and class estimates, as in
Theorem~\ref{thm:mb-branch-sup}; no weak-limit transfer of a
supremum is invoked.  Physical reversal preserves arclength and
source-angle measure and reverses momentum signs.  It proves the
backward assertions without changing the conditioning event.
\end{proof}

\begin{corollary}[The selected angular source mass]\label{cor:lc-source-mass}
For either source after its first true collision, and every $h\ge0$,
\begin{equation}\label{eq:lc-source-mass}
 \lambda\{T_R^hY\in E_{\delta,R}\}
                      \le A_{\rm src}G_\ell L_{\rm c}\delta.
\end{equation}
For the horizontal source selected before its first collision,
\begin{equation}\label{eq:lc-horizontal-initial}
 \int_{F_a}\one_{E_{\delta,R}}\,d\alpha_0
                       \le J_0L_{\rm c}\delta/R_-.
\end{equation}
These are one-dimensional physical source estimates.  They will
become cap-boundary charges only after the exact Jacobian
cancellation proved in the next section.
\end{corollary}
\begin{proof}
The total transported density supremum is at most
$A_{\rm src}G_\ell$.  Integrate it on the clearance covers of
combined length at most $L_{\rm c}\delta$, and sum the exact
labelled sources.  For the initial horizontal fibre divide the
circle into the $J_0$ preliminary arcs and apply the horizontal
cover directly.  Since $ds=R\,d\alpha_0$, the conversion costs
at most $R_-^{-1}$.  Every retained original observation is restored
by positive restriction.  No selected event has been replaced by
a normalized fibre law.
\end{proof}
